\scanpage{22}

\mspage{10}

sûrement vrai que $SF_M$ est encore connexe, mais je ne veux pas m'y attarder maintenant.

Je reviens au cas général \deleted{\unsure{(}} (avec tjrs $2.1_\Lambda\ne0$, cependant). On a donc

\[
\mathrm{SL}(M)'\overset{\mathrm{déf}}{=}\mathrm{SL}(M)/\{\pm1\}\longrightarrow
\underbrace{\operatorname{Aut}^{+}(SP_M)\deleted{\unsure{\{\mu\}}}}_{\substack{\text{pour l'orientation}\\\text{canonique }\mu.}}\tag*{(48)}
\]

Je dis que cette application est injective, \deleted{Soit en effet $u\in$} \deleted{i.e. que} plus précisément

\[
\left\{\begin{aligned}
&\text{Soit } f\in\deleted{S}F_M,\quad\text{et } u\in\mathrm{GL}(M),\quad\text{alors}\\
&(u\deleted{\unsure{(}}(x)=x\quad\forall x\in f)\Longrightarrow u\in\{\pm\mathrm{id}_M\}.
\end{aligned}\right.\tag*{(49)}
\]

OPS $f=\{\dot{x},\dot{y},\dot{z}\}$, avec $(x,y,z)\in B(M)$, on aura donc, quitte à changer $u$ de signe

\[
ux=x,\quad uy=\varepsilon'y,\quad uz=\varepsilon''z\qquad\deleted{x+y+}
\]

\[
\text{\unsure{avec}}\quad\deleted{\unsure{z=}}\quad x+y+z=0,\quad\underbrace{x+\varepsilon'y+\varepsilon''z}_{=u(x+y+z)}=0
\]

\[
\text{donc}\quad \deleted{z=}\quad x=\deleted{\unsure{\varepsilon}}-y-z=-\varepsilon'y-\varepsilon''z
\]

et comme $(y,z)\in\operatorname{Rep}(M)$, cel implique

\[
\varepsilon'=1,\ \varepsilon''=1\qquad\text{donc}\qquad u=\mathrm{id}_M\qquad\text{OK.}
\]

Ainsi, on trouve un homomorphisme injectif de suites exactes canoniques

% archive sheet number: 21
